% Generated draft: requires mathematical and source review.
\documentclass[11pt]{article}
\usepackage[margin=25mm]{geometry}
\usepackage{fontspec}
\setmainfont{Latin Modern Roman}
\usepackage{amsmath,amssymb,mathtools,unicode-math}
\setmathfont{Latin Modern Math}
\usepackage{enumitem,needspace,xurl,hyperref}
\setlength{\parindent}{0pt}
\setlength{\parskip}{5pt}
\setlength{\emergencystretch}{3em}
\newcommand{\N}{\mathbb{N}}
\newcommand{\Z}{\mathbb{Z}}
\newcommand{\Q}{\mathbb{Q}}
\newcommand{\R}{\mathbb{R}}
\newcommand{\C}{\mathbb{C}}
\newcommand{\F}{\mathbb{F}}
\newcommand{\A}{\mathbb{A}}
\newcommand{\PP}{\mathbb{P}}
\newcommand{\E}{\mathbb{E}}
\newcommand{\eps}{\varepsilon}
\newcommand{\dd}{\,\mathrm{d}}
\begin{document}
\section*{UnsolvedMath ID: 30005886}
\textbf{Source problem code:} catalogue entry
\subsection{Definitions and mathematical statement}
For each integer \(n\geq 3\), let \(\operatorname{GL}_n(\Z)\) denote the group of integral \(n\times n\) matrices with determinant \(\pm1\). For an ordered pair \((A,B)\in\operatorname{GL}_n(\Z)^2\), put
\[
\Gamma=\langle A,B\rangle
=\{X_1\cdots X_k:k\geq0,\ X_i\in\{A,A^{-1},B,B^{-1}\}\},
\]
where the empty product is the identity. Let \(G\) be the Zariski closure of \(\Gamma\) in the algebraic group \(\operatorname{GL}_{n,\Q}\), and use the given matrix embedding to define
\[
G(\Z)=G(\Q)\cap\operatorname{GL}_n(\Z).
\]
Using the usual infinite-index convention, \(\Gamma\) is \emph{thin in its Zariski closure} if it is Zariski dense in \(G\) and
\[
[G(\Z):\Gamma]=\infty.
\]
The density condition holds by the definition of \(G\).

\textbf{Classification problem.} For every integer \(n\geq3\), classify all pairs
\[
(A,B)\in\operatorname{GL}_n(\Z)^2
\quad\text{such that}\quad
[G(\Z):\langle A,B\rangle]=\infty,
\qquad G=\overline{\langle A,B\rangle}^{\,\mathrm{Zar}}.
\]
No assumptions of connectedness, semisimplicity, irreducibility, or density in the full general linear group are supplied. The source does not specify an equivalence relation on pairs or the desired form of a classification. The interpretation of \(G\) as the Zariski closure follows the supplied clean statement; the original excerpt writes \(G(\Z)\) without defining \(G\). The definition above makes explicit the conventional meaning of thinness used in this draft, which requires source review.
\subsection{Short English statement}
For every integer \(n\geq3\), determine which pairs of invertible integral matrices generate a subgroup of infinite index in the integral points of its own Zariski closure. The intended equivalence relation for the classification is unspecified. The supplied record labels the problem open; this draft does not independently verify that status.
\subsection{Sources}
\begin{itemize}
\item Supplied local UnsolvedMath registry record, permanent ID \(30005886\), title \emph{Classifying Thin Two-Generator Matrix Groups}, exact source code \texttt{catalogue entry}, in the Oberwolfach Reports Open Problems collection.
\item Supplied report citation: \emph{Mini-Workshop: Growth and Expansion in Groups} (2024), Oberwolfach Reports, pp.\ 1005--1040, DOI \texttt{10.4171/owr/2024/17}. \url{https://doi.org/10.4171/owr/2024/17}. The original excerpt attributes the questions to S.\ Eberhard; the separate proposer field is empty.
\item Supplied literature source labeled ``generic thinness'': \url{https://academic.oup.com/imrn/article/2017/17/5385/3056825}.
\item Supplied literature source labeled ``thin monodromy classification'': \url{https://real.mtak.hu/212740/}.
\item Supplied literature source labeled ``OWR report'': \url{https://doi.org/10.4171/owr/2024/17}.
\item Literature assessment dated exactly \texttt{2026-08-22}: the supplied curation labels the problem open and reports that special families, notably hypergeometric monodromy, have been classified, that generic thinness is known, and that no classification of all such integral matrix pairs is available. These are supplied dated curation claims, not independently verified findings. The supplied background explicitly requests independent verification of time-sensitive status and cited claims.
\item The original excerpt also appends a separate question about the shortest law on \(S_n\), attributing an upper bound to Kozma--Thom, reference [15]. No full bibliographic entry for [15] is supplied.
\end{itemize}
\end{document}
